
\documentclass[11pt,a4paper]{article}

\usepackage[margin=2.5cm]{geometry}
\usepackage{amsmath}
\usepackage{hyperref}
\usepackage{booktabs}
\usepackage{enumitem}

\title{Initial Project Documentation\\
Deep Learning Models Comparison}
\author{}
\date{\today}

\begin{document}

\maketitle

\section{Project Overview}

The goal of this project is to experimentally evaluate and compare different deep learning approaches for image classification. The study focuses on:

\begin{itemize}
    \item convolutional neural networks (CNN),
    \item few-shot learning methods,
    \item modern pretrained architectures,
    \item ensemble techniques.
\end{itemize}

The project investigates how architectural choices, training hyperparameters, regularization strategies, and data augmentation techniques influence model performance.

All experiments will be conducted using PyTorch with training performed locally on GPU hardware. Experiment tracking and visualization will be handled using TensorBoard.

\section{Technology Stack}

The project will use the following tools:

\begin{itemize}
    \item framework: PyTorch;
    \item hardware: local GPU;
    \item experiment tracking: TensorBoard;
    \item programming language: Python;
    \item version control: Git.
\end{itemize}

PyTorch was selected due to its flexibility and strong support for experimental workflows.

\section{Data Augmentation}

To improve generalization and analyze augmentation influence, augmented datasets will be generated and stored prior to training.

\subsection{Standard Augmentation Techniques}

The following standard transformations will be applied:

\begin{itemize}
    \item horizontal flip,
    \item rotation by $45^\circ$ (random direction),
    \item Gaussian blur.
\end{itemize}

\subsection{Advanced Augmentation Techniques}

Four advanced augmentation methods will be used:

\begin{itemize}
    \item Cutout - randomly masking square regions of the image,
    \item color jittering - transforming saturation, brightness and hue,
    \item Autoaugment - python library for automatic and optimal data augmentation,
    \item creating noise and artifacts using image compression.
\end{itemize}


\section{Experimental Methodology}

The experimental pipeline consists of the following stages:

\begin{enumerate}
    \item Training and analysis of multiple CNN configurations.
    \item Selection of the best-performing CNN (CNN representative).
    \item Comparison with:
    \begin{itemize}
        \item few-shot learning model,
        \item EfficientNet architecture.
    \end{itemize}
    \item Ensemble learning using soft voting.
\end{enumerate}

\section{CNN Implementation and Experiments}

The CNN architecture will follow the official PyTorch tutorial:

\begin{center}
\url{https://docs.pytorch.org/tutorials/beginner/blitz/cifar10_tutorial.html}
\end{center}

This architecture serves as a controlled baseline for systematic hyperparameter analysis.

\subsection{Training Hyperparameters}

The following training-related hyperparameters will be investigated:

\begin{itemize}
    \item Optimizer:
    \begin{itemize}
        \item SGD
        \item Adam
    \end{itemize}
    \item Momentum (SGD optimizer, 2-3 different values)
\end{itemize}

\subsection{Regularization Hyperparameters}

The following regularization parameters will be analyzed:

\begin{itemize}
    \item Dropout:
    \begin{itemize}
        \item disabled
        \item enabled
    \end{itemize}
    \item Cross-Entropy label smoothing:
    \begin{itemize}
        \item zero smoothing
        \item non-zero smoothing
    \end{itemize}
\end{itemize}

\subsection{Experimental Grid}

Each CNN variant will be evaluated across combinations of:

\begin{itemize}
    \item augmentation strategy (including selected combinations of more than one),
    \item optimizer type,
    \item momentum value,
    \item label smoothing,
    \item dropout configuration.
\end{itemize}

The best-performing configuration will be selected as the \textbf{CNN representative model}.

\section{Few-Shot Learning}

To address learning under limited data conditions, a few-shot learning method will be implemented.

\subsection{Method Used}

Prototypical Networks implementation:

\begin{center}
\url{https://github.com/orobix/Prototypical-Networks-for-Few-shot-Learning-PyTorch}
\end{center}

Prototypical networks learn a metric space where classification is performed using distances to class prototypes.

\section{EfficientNet Comparison}

A pretrained EfficientNet model will be used as a modern architecture baseline.

EfficientNet provides:

\begin{itemize}
    \item compound scaling of depth, width, and resolution,
    \item strong performance with relatively low parameter count,
    \item transfer learning capabilities.
\end{itemize}

The network will be fine-tuned on the project dataset and compared with other models.

\section{Reduced Training Dataset Experiments}

To analyze robustness to limited data availability:

\begin{enumerate}
    \item A reduced subset of the training dataset will be created.
    \item Selected models will be trained using:
    \begin{itemize}
        \item full dataset,
        \item reduced dataset.
    \end{itemize}
    \item Performance degradation will be analyzed.
\end{enumerate}

\section{Ensemble Method --- Soft Voting}

An ensemble model will be implemented using soft voting.

Each model outputs class probability distributions:

\begin{itemize}
    \item CNN representative,
    \item few-shot model,
    \item EfficientNet.
\end{itemize}

Final prediction is computed as:

\begin{equation}
P_{\text{final}} = \frac{1}{3}\sum_{i=1}^{3} P_i
\end{equation}

where $P_i$ denotes predicted probabilities from model $i$.

\section{Evaluation Metrics}

Models will be evaluated using:

\begin{itemize}
    \item classification accuracy
    \item validation loss
    \item precision
    \item recall
\end{itemize}

Training dynamics will be monitored using TensorBoard.

\section{Experimental Procedure}

Each experiment follows:

\begin{enumerate}
    \item Dataset preparation and augmentation
    \item Model initialization
    \item Training under fixed epoch budget
    \item Validation evaluation
    \item Logging results in TensorBoard
    \item Saving best checkpoints
    \item Comparative analysis
\end{enumerate}

To ensure fairness:

\begin{itemize}
    \item identical train/validation splits size will be used,
    \item random seeds will be fixed where possible.
\end{itemize}

\section{Reproducibility}

Reproducibility will be ensured by:

\begin{itemize}
    \item fixing random seeds,
    \item storing augmented datasets,
    \item logging hyperparameters,
    \item saving trained checkpoints,
    \item version-controlling the source code.
\end{itemize}

\end{document}
